% GENERATED FILE. Edit canonical lessons, then run npm run generate:books.
% canonical-source-hash: fnv1a64:7ce5a17159db86e8
% canonical-lessons: LA-C373-effectus, LA-C373-consequor, LA-C373-provenio, LA-C373-succedo, LA-C373-redundo

\chapter{How Things Turn Out}
\label{ch:la-how-things-turn-out}
\hlchaptermodality{373}

\begin{chapteropening}
\textbf{By the end of this chapter:} \emph{I can speak of a result, say what someone achieved and how things turned out, and that blame or credit falls on someone.}

Complete the last lesson of chapter 373: I can speak of a result, say what someone achieved and how things turned out, and that blame or credit falls on someone.
\end{chapteropening}

\section[effectus, effectūs]{effectus, effectūs — a result, an effect; performance}

\label{lesson:LA-C373-effectus}



\begin{quote}
Before the new one: say the Latin for effective, then the Latin for to rouse.
\end{quote}

\subsection*{You'll want to know: effectus, effectūs}
\textbf{effectus, effectūs} — \textquotedblleft{}a result, an effect; performance\textquotedblright{}.

\begin{quote}
Effectus cōnsiliī bonus fuit.
\end{quote}

\begin{quote}
\emph{The result of the plan was good.}
\end{quote}

\emph{effectus, effectūs} m., built on \emph{efficiō}. English \emph{effect}.

\begin{cousinweb}
\emph{effectus} is the noun built on \emph{efficiō} \textquotedblleft{}to bring about\textquotedblright{} (EtymDB; Lewis \& Short). English \emph{effect} comes from it.
\end{cousinweb}

\begin{grammarlens}[title={What you've built}]
One more word for telling, explaining and describing at B1.
\end{grammarlens}

\subsection*{Your turn}
\begin{itemize}
\raggedright
  \item \emph{Say it:} \emph{effectus, effectūs}
  \item \emph{Say it:} \emph{effectus, effectūs}, once more
  \item \emph{Say it:} a sentence of your own with \emph{effectus, effectūs}
  \item \emph{Recall:} say the Latin for effective, then the Latin for to rouse, then say \emph{effectus, effectūs} again
\end{itemize}

\begin{tcolorbox}[breakable,colback=teal!4,colframe=teal!35!black,title={Before you move on}]
What does effectus, effectūs mean? (\textquotedblleft{}A result, an effect; performance\textquotedblright{}.) Say it once more.
\end{tcolorbox}
\section[cōnsequor, cōnsequī]{cōnsequor, cōnsequī — to follow; to achieve, to obtain}

\label{lesson:LA-C373-consequor}



\begin{quote}
Before the new one: say the Latin for to rouse, then the Latin for a result.
\end{quote}

\subsection*{You'll want to know: cōnsequor, cōnsequī}
\textbf{cōnsequor, cōnsequī} — \textquotedblleft{}to follow; to achieve, to obtain\textquotedblright{}.

\begin{quote}
Quod voluit, cōnsecūtus est.
\end{quote}

\begin{quote}
\emph{He achieved what he wanted.}
\end{quote}

Deponent: \emph{cōnsequor, cōnsequī, cōnsecūtus sum}. \emph{Cōnsequitur ut}, it follows that. English \emph{consequence}.

\begin{cousinweb}
\emph{cōnsequor} is built from \emph{con-} and \emph{sequor} \textquotedblleft{}to follow\textquotedblright{} (EtymDB; Lewis \& Short): to follow all the way to the goal. English \emph{consequence} comes from its family.
\end{cousinweb}

\begin{grammarlens}[title={What you've built}]
One more word for telling, explaining and describing at B1.
\end{grammarlens}

\subsection*{Your turn}
\begin{itemize}
\raggedright
  \item \emph{Say it:} \emph{cōnsequor, cōnsequī}
  \item \emph{Say it:} \emph{cōnsequor, cōnsequī}, once more
  \item \emph{Say it:} a sentence of your own with \emph{cōnsequor, cōnsequī}
  \item \emph{Recall:} say the Latin for to rouse, then the Latin for a result, then say \emph{cōnsequor, cōnsequī} again
\end{itemize}

\begin{tcolorbox}[breakable,colback=teal!4,colframe=teal!35!black,title={Before you move on}]
What does cōnsequor, cōnsequī mean? (\textquotedblleft{}To follow; to achieve, to obtain\textquotedblright{}.) Say it once more.
\end{tcolorbox}
\section[prōveniō, prōvenīre]{prōveniō, prōvenīre — to come forth, to grow; to turn out}

\label{lesson:LA-C373-provenio}



\begin{quote}
Before the new one: say the Latin for a result, then the Latin for to follow.
\end{quote}

\subsection*{You'll want to know: prōveniō, prōvenīre}
\textbf{prōveniō, prōvenīre} — \textquotedblleft{}to come forth, to grow; to turn out\textquotedblright{}.

\begin{quote}
Frūmentum hōc annō bene prōvēnit.
\end{quote}

\begin{quote}
\emph{The grain grew well this year.}
\end{quote}

\emph{prōveniō, prōvenīre, prōvēnī, prōventum}. Of a harvest, to come up; of a plan, to turn out. English \emph{provenance} comes from it through French.

\begin{cousinweb}
\emph{prōveniō} is built from \emph{prō} \textquotedblleft{}forth\textquotedblright{} and \emph{veniō} \textquotedblleft{}to come\textquotedblright{} (EtymDB; Lewis \& Short). English \emph{provenance} comes from it through French.
\end{cousinweb}

\begin{grammarlens}[title={What you've built}]
One more word for telling, explaining and describing at B1.
\end{grammarlens}

\subsection*{Your turn}
\begin{itemize}
\raggedright
  \item \emph{Say it:} \emph{prōveniō, prōvenīre}
  \item \emph{Say it:} \emph{prōveniō, prōvenīre}, once more
  \item \emph{Say it:} a sentence of your own with \emph{prōveniō, prōvenīre}
  \item \emph{Recall:} say the Latin for a result, then the Latin for to follow, then say \emph{prōveniō, prōvenīre} again
\end{itemize}

\begin{tcolorbox}[breakable,colback=teal!4,colframe=teal!35!black,title={Before you move on}]
What does prōveniō, prōvenīre mean? (\textquotedblleft{}To come forth, to grow; to turn out\textquotedblright{}.) Say it once more.
\end{tcolorbox}
\section[succēdō, succēdere]{succēdō, succēdere — to turn out well, to succeed; to come next}

\label{lesson:LA-C373-succedo}



\begin{quote}
Before the new one: say the Latin for to follow, then the Latin for to come forth.
\end{quote}

\subsection*{You'll want to know: succēdō, succēdere}
\textbf{succēdō, succēdere} — \textquotedblleft{}to turn out well, to succeed; to come next\textquotedblright{}.

\begin{quote}
Cōnsilium nōn successit.
\end{quote}

\begin{quote}
\emph{The plan did not succeed.}
\end{quote}

\emph{succēdō, succēdere, successī, successum}. With the dative it is also to take someone's place. English \emph{succeed} and \emph{success}.

\begin{cousinweb}
\emph{succēdō} is built from \emph{sub} \textquotedblleft{}up to, next\textquotedblright{} and \emph{cēdō} \textquotedblleft{}to go\textquotedblright{} (EtymDB; Lewis \& Short): to come up close, to come next. English \emph{succeed} and \emph{success} come from it.
\end{cousinweb}

\begin{grammarlens}[title={What you've built}]
One more word for telling, explaining and describing at B1.
\end{grammarlens}

\subsection*{Your turn}
\begin{itemize}
\raggedright
  \item \emph{Say it:} \emph{succēdō, succēdere}
  \item \emph{Say it:} \emph{succēdō, succēdere}, once more
  \item \emph{Say it:} a sentence of your own with \emph{succēdō, succēdere}
  \item \emph{Recall:} say the Latin for to follow, then the Latin for to come forth, then say \emph{succēdō, succēdere} again
\end{itemize}

\begin{tcolorbox}[breakable,colback=teal!4,colframe=teal!35!black,title={Before you move on}]
What does succēdō, succēdere mean? (\textquotedblleft{}To turn out well, to succeed; to come next\textquotedblright{}.) Say it once more.
\end{tcolorbox}
\section[redundō, redundāre]{redundō, redundāre — to overflow; to result (in), to fall back (on)}

\label{lesson:LA-C373-redundo}



\begin{quote}
Before the new one: say the Latin for to come forth, then the Latin for to turn out well.
\end{quote}

\subsection*{You'll want to know: redundō, redundāre}
\textbf{redundō, redundāre} — \textquotedblleft{}to overflow; to result (in), to fall back (on)\textquotedblright{}.

\begin{quote}
Haec culpa in tē redundat.
\end{quote}

\begin{quote}
\emph{This fault falls back on you.}
\end{quote}

\emph{redundō, redundāre, redundāvī, redundātum}. Built on \emph{unda}, a wave: water flowing back over, and blame or credit coming back. English \emph{redound} and \emph{redundant}.

\begin{cousinweb}
\emph{redundō} is built from \emph{red-} \textquotedblleft{}back\textquotedblright{} and \emph{undō} \textquotedblleft{}to surge\textquotedblright{}, from \emph{unda} \textquotedblleft{}a wave\textquotedblright{} (Lewis \& Short). English \emph{redound} and \emph{redundant} come from it.
\end{cousinweb}

\begin{grammarlens}[title={What you've built}]
One more word for telling, explaining and describing at B1.
\end{grammarlens}

\subsection*{Your turn}
\begin{itemize}
\raggedright
  \item \emph{Say it:} \emph{redundō, redundāre}
  \item \emph{Say it:} \emph{redundō, redundāre}, once more
  \item \emph{Say it:} a sentence of your own with \emph{redundō, redundāre}
  \item \emph{Recall:} say the Latin for to come forth, then the Latin for to turn out well, then say \emph{redundō, redundāre} again
\end{itemize}

\begin{tcolorbox}[breakable,colback=teal!4,colframe=teal!35!black,title={Before you move on}]
What does redundō, redundāre mean? (\textquotedblleft{}To overflow; to result (in), to fall back (on)\textquotedblright{}.) Say it once more.
\end{tcolorbox}
